\documentclass[12pt,a4paper]{article}
\usepackage[UTF8,heading=true]{ctex}
\usepackage{geometry}
\usepackage{amsmath,amssymb,amsthm}
\usepackage{graphicx}
\usepackage{booktabs}
\usepackage{longtable}
\usepackage{array}
\usepackage{enumitem}
\usepackage{hyperref}
\usepackage{fancyhdr}
\usepackage{xcolor}
\usepackage{listings}
\usepackage{setspace}  % 行距控制包

% 配置 ctex 章节标题为 "一、"
\ctexset{
  section = {
    name = {,、},
    number = \chinese{section},
    format = {\Large\bfseries\raggedright}
  }
}

\geometry{left=2.5cm,right=2.5cm,top=2.5cm,bottom=2.5cm}
\setlength{\headheight}{15pt}
\setstretch{1.5}  % 设置1.5倍行距（CNIPA常用）

\pagestyle{fancy}
\fancyhf{}
\fancyhead[C]{发明专利技术交底书}
\fancyfoot[C]{\thepage}

\title{\textbf{发明专利技术交底书} \\[0.5em]
\Large 一种具备内生门控计算的混合专家模型片上网络路由器及方法}
\author{中国移动通信有限公司研究院}
\date{\today}

\begin{document}

\begin{center}
    {\Large\heiti 中国移动专利申请} \\[1em]
    {\Huge\heiti 技术交底书} \\[1em]
\end{center}

\noindent
\begin{tabular}{|p{3cm}|p{12cm}|}
\hline
公司编号 &  \\
\hline
发明名称 & 一种具备内生门控计算的混合专家模型片上网络路由器及方法 \\
\hline
申报单位 & 研究院 \\
\hline
申报类型 & 发明 \\
\hline
发明人 & 许达 \\
\hline
技术联系人 & 许达 (xudayj@chinamobile.com, +86-13521894156) \\
\hline
注意事项 & 1．技术联系人应为深入了解本申请提案技术方案的技术人员。 \\
         & 2．请按照集团公司提供的本技术交底书模板逐项填写。 \\
\hline
\end{tabular}

\vspace{1cm}

%==============================================================================
\section{发明名称}
%==============================================================================

一种具备内生门控计算的混合专家模型片上网络路由器及方法

%==============================================================================
\section{技术领域}
%==============================================================================

本发明涉及人工智能推理加速芯片、片上网络（Network-on-Chip, NoC）、并行计算体系结构与软硬协同优化领域，尤其涉及一种面向混合专家模型（Mixture-of-Experts, MoE）条件路由计算的片上网络路由器结构、路由与预取控制方法，以及与编译器/运行时协同的门控（Gating）计算下沉机制。

\textbf{与量子计算的关联说明：}

当前大模型推理正从“纯稠密矩阵乘法”演化为“稠密计算 + 动态条件路由”的组合形态。MoE 模型通过门控网络对每个 token 动态选择少量专家，从而在参数规模很大时仍保持较低的计算开销。然而，MoE 的关键瓶颈不再是乘加器数量，而是跨专家的数据搬运、路由等待、以及由软件调度引入的流水线空泡（Pipeline Stall）与片上网络拥塞。

本发明提出 Algorithm-to-Silicon 的思路：将 MoE 的门控路由决策与数据包路由控制在片上网络路由器内部“同周期融合”执行，降低软件参与度与等待开销，提升 MoE 推理吞吐与能效。

\textbf{关联项目：}\underline{\hspace{8cm}}（待填写）

%==============================================================================
\section{现有技术的技术方案}
%==============================================================================

\subsection{量子计算网络安全背景}

MoE 模型通常在每层采用一个门控网络对输入 token 计算专家得分，选择 Top-$k$ 专家执行后续计算。现有系统往往在计算核心（如 GPU/TPU/NPU）上完成门控得分计算，并由软件运行时或编译器在门控结果产生后安排数据搬运与专家执行。该流程存在以下特征：
\begin{enumerate}
    \item 门控结果产生与数据搬运存在严格的时序依赖，导致核心等待；
    \item token-to-expert 的稀疏路由呈现强动态性与突发性，使得 NoC 容易出现热点拥塞；
    \item 门控与路由的实现通常依赖软件队列、DMA 触发与同步屏障，带来额外的控制开销。
\end{enumerate}

\subsection{后量子密码技术现状}

现有 NoC 路由器主要面向“已知目的地”的确定性/自适应路由：在数据包头部携带目的地坐标，路由器按路由算法选择输出端口，辅以虚通道（VC）与信用（credit）流控。针对 AI 加速器的改进多集中在带宽与拥塞管理，例如增加多播/归约支持、改进仲裁策略或配置更大缓冲。

对于 MoE 的条件路由，现有方案普遍仍采用“先算门控、后路由搬运”的解耦方式：门控由计算核心完成，路由器只负责搬运，无法在数据经过 NoC 的过程中主动完成门控决策或预测性预取。因此，当门控计算结果出来时，往往已经错过了可用于隐藏通信延迟的时间窗口。

%==============================================================================
\section{现有技术的缺点及本申请提案要解决的技术问题}
%==============================================================================

现有技术在 MoE 推理场景下存在如下缺点：
\begin{enumerate}
    \item \textbf{门控与通信割裂导致等待}：门控结果产生后才触发搬运，形成串行依赖，造成流水线空泡；
    \item \textbf{软件调度开销高}：运行时需要维护 token 队列、排序/分桶、触发 DMA、同步屏障等，消耗主核/控制核资源；
    \item \textbf{片上网络拥塞难以预防}：token 的专家选择呈现稀疏但突发的流量模式，导致瞬时热点、队头阻塞（HoL blocking）；
    \item \textbf{HBM/SRAM 预取缺乏前瞻}：专家权重或激活在被真正需要时才搬运，无法充分利用 NoC 空隙进行异步预取。
\end{enumerate}

因此，本申请提案要解决的技术问题是：在不依赖新材料、仅使用成熟 CMOS 数字逻辑的前提下，提供一种 NoC 路由器微架构与方法，使其能够在转发 token 数据包的同时内生执行门控相关的轻量计算，完成专家选择、预测路由与 DMA 前瞻预取，从而降低软件参与、减少流水线等待并提升 MoE 推理整体吞吐与能效。

%==============================================================================
\section{本申请提案的技术方案的详细阐述}
%==============================================================================

\subsection{系统架构}

本发明提出一种“具备内生门控计算（In-Network Gating）的 MoE 片上网络路由器”（以下简称路由器或 InNG-Router），其可集成于 AI 推理加速芯片的 NoC 交换节点中，与计算核心、片上 SRAM/HBM 控制器、DMA 引擎以及编译器/运行时协同工作。其总体组成包括：
\begin{enumerate}
    \item \textbf{门控特征提取与量化单元（Gating Feature Unit, GFU）}：从 token 数据包头或侧带元数据中提取门控输入特征（如门控向量索引、低维投影、或压缩 embedding），并执行低比特量化/截断；
    \item \textbf{轻量点积计算阵列（Micro-Dot Engine, MDE）}：在路由器内置少量乘加逻辑，对门控特征与专家路由权重（或其压缩版本）计算 Top-$k$ 近似得分；
    \item \textbf{预测路由与拥塞感知仲裁器（Predictive Routing Arbiter, PRA）}：根据得分与拥塞状态选择输出端口/虚通道，支持 lookahead 预测与多目标复制（当 $k>1$）；
    \item \textbf{前瞻 DMA 预取触发器（Lookahead Prefetch Trigger, LPT）}：在 token 尚未到达目标专家核心前，向 DMA/存储控制器发起专家权重或上下文块的预取请求；
    \item \textbf{一致性与重排序支持}：为 token 的多专家并行与回写合并提供序号标签（token id/layer id）与轻量重排序缓冲。
\end{enumerate}

为保证通用性，门控计算可配置为以下两种模式之一：
\begin{itemize}
    \item \textbf{模式A（下沉门控）}：门控网络的关键计算（或其近似）在路由器内执行，计算核心只产生门控输入特征；
    \item \textbf{模式B（辅助门控）}：门控得分已由计算核心产生，路由器仅执行预测路由、拥塞仲裁与预取触发。
\end{itemize}

\subsection{核心方法步骤}

\subsubsection{1. 分布式密钥生成 (KeyGen)}

编译器/运行时在生成 MoE 层的执行计划时，将每个 token 的门控输入特征编码为路由器可识别的侧带字段，随 token 数据包一并注入 NoC。为降低开销，特征可采用以下之一：
\begin{enumerate}
    \item 低维投影后的向量索引或码本编号；
    \item 低比特量化的门控向量片段（如 INT8/INT4）；
    \item 专家选择的候选集合提示（候选掩码或哈希）。
\end{enumerate}
路由器的 GFU 对该字段解码并完成必要的对齐与量化处理，输出供 MDE 使用的门控特征表示。

\subsubsection{2. 消息哈希与映射}

路由器 MDE 存储或缓存一份用于专家打分的路由权重（可为压缩/分段形式）。在 token 数据包经过路由器的同一调度窗口内，MDE 对门控特征执行近似点积计算：
\[
s_e = \langle g, r_e \rangle,\quad e \in \{1,\dots,E\}
\]
其中 $g$ 为门控特征，$r_e$ 为第 $e$ 个专家的路由权重（或其近似表示）。为降低逻辑规模，可采用分组候选、局部 Top-$m$、或两级选择（粗筛 + 精筛）方式实现 Top-$k$ 专家选择。PRA 将 Top-$k$ 结果编码为目的端口/目的坐标集合，写入数据包头部或路由器内部转发表项。

\subsubsection{3. 分布式高斯采样}

为降低拥塞与队头阻塞，PRA 结合当前路由器缓冲占用、下游 credit、以及历史流量统计，对输出端口与虚通道进行预测性选择。其可采用：
\begin{enumerate}
    \item \textbf{lookahead}：预估后续若干跳的拥塞概率，提前规避热点路径；
    \item \textbf{多目标复制与限流}：当 $k>1$ 需复制 token 至多个专家时，按拥塞分配复制时序并限制并发复制数；
    \item \textbf{快速回退}：若下游拥塞突发，允许在不改变专家集合的前提下选择替代路径或延后发送。
\end{enumerate}
该步骤在路由器仲裁阶段完成，不依赖软件同步。

\subsubsection{4. 基于 NTT 的线性变换}

当路由器确定 token 的目标专家集合后，LPT 在 token 尚未到达专家核心前，向 DMA/存储控制器发起预取请求。预取对象可包括：
\begin{enumerate}
    \item 专家权重分块（按层/按专家/按 tile 切分）；
    \item 专家侧激活缓冲或中间缓存行；
    \item 需要参与合并的输出缓冲段。
\end{enumerate}
为避免无效预取，LPT 可结合门控得分阈值、token 批次一致性与历史命中率，动态调整预取粒度与并发度，并对预取请求进行去重与合并。

\subsubsection{5. 基于 Beaver 三元组的安全范数验证}

针对 MoE 的“分发-并行-合并”模式，路由器为每个 token 附加序号标签（token id, layer id, expert id），并在必要时维护轻量重排序缓冲，保证：
\begin{enumerate}
    \item 多专家输出在合并点按照 token 序一致地归并；
    \item 拥塞导致的乱序到达不会破坏上层计算语义；
    \item credit/VC 流控与预取触发不会产生死锁。
\end{enumerate}
合并控制可由下游路由器或专用合并节点完成，路由器侧仅需提供标签与最小缓冲支持。

\subsubsection{6. 动态节点管理}

为支持专家在芯片上的动态放置（如不同批次/不同层的专家映射变化），本发明提供路由表更新机制：编译器/运行时在层切换或负载迁移时，向 NoC 路由器下发专家映射表的增量更新。路由器可在不停止数据面转发的情况下更新控制面表项，并支持双缓冲切换以保证一致性。

%==============================================================================
\section{本申请提案的关键点和欲保护点}
%==============================================================================

本申请提案的关键点和欲保护点包括但不限于：
\begin{enumerate}
    \item 一种片上网络路由器，其在转发 token 数据包的同时内生执行门控相关的轻量点积计算，并输出 Top-$k$ 专家选择结果，用于驱动数据包目的路由；
    \item 一种路由器内置的预测路由与拥塞感知仲裁机制，用于在 MoE 动态稀疏流量下减少热点拥塞与队头阻塞；
    \item 一种由路由器根据门控结果触发的 DMA 前瞻预取机制，实现专家权重/上下文块在 token 到达前的异步预取与去重合并；
    \item 一种端到端的标签化与轻量重排序机制，支撑 token 的多专家并行分发与合并语义；
    \item 一种专家映射表的可重构更新方法，使得专家在芯片上的放置变化时路由器可在线更新而不影响数据面。
\end{enumerate}

%==============================================================================
\section{与第三条中最接近的现有技术相比，本申请提案有何技术优点}
%==============================================================================

与现有“门控在计算核心、路由在 NoC”的解耦方案相比，本申请提案具有显著技术优点：
\begin{enumerate}
    \item \textbf{降低流水线空泡}：门控决策与路由转发同窗口执行，减少“算完门控再等待搬运”的串行依赖；
    \item \textbf{降低软件控制开销}：减少运行时的分桶、调度、同步屏障与 DMA 触发次数，实现更接近硬件流水的 MoE 执行；
    \item \textbf{提升通信可预测性}：拥塞感知与预测路由降低热点概率，使 token-to-expert 流量更平滑；
    \item \textbf{隐藏存储访问延迟}：路由器前瞻预取在数据到达前准备权重/缓存块，提高 SRAM/HBM 命中，提升吞吐与能效；
    \item \textbf{工程落地简单}：仅需在路由器控制/仲裁逻辑中增加小规模点积阵列与控制面扩展，不依赖新材料或复杂工艺。
\end{enumerate}

%==============================================================================
\section{其他有助于理解本申请提案的技术资料}
%==============================================================================

(1) Shazeer, N., et al. ``Outrageously Large Neural Networks: The Sparsely-Gated Mixture-of-Experts Layer.'' ICLR, 2017. \\
(2) Fedus, W., Zoph, B., Shazeer, N. ``Switch Transformers: Scaling to Trillion Parameter Models with Simple and Efficient Sparsity.'' JMLR/ArXiv, 2021. \\
(3) Lepikhin, D., et al. ``GShard: Scaling Giant Models with Conditional Computation and Automatic Sharding.'' ICML, 2021. \\
(4) Dally, W.J., Towles, B. ``Principles and Practices of Interconnection Networks.'' Morgan Kaufmann, 2004. \\
(5) Kim, J., et al. ``A Lightweight Router for Network-on-Chip.'' (相关 NoC 路由/仲裁研究综述文献，可据具体实现补充).

%==============================================================================
\section{本申请提案的侵权证据可获得性}
%==============================================================================

本申请提案的侵权证据可获得性较高，主要体现在：
\begin{enumerate}
    \item \textbf{微架构可观测特征}：芯片公开资料、白皮书或专利中常会披露 NoC 路由器是否具备 in-network compute、门控下沉、预测路由与预取触发等特性；
    \item \textbf{性能行为特征}：在公开基准（如 MoE 推理吞吐、跨专家通信开销、stall 比例）中，若出现与本发明一致的“门控与通信融合带来的显著 stall 降低”，可作为线索；
    \item \textbf{软件接口与编译器痕迹}：若对外暴露 MoE 路由元数据格式、专家映射表更新接口、或预取提示指令/ABI，容易在编译器产物与运行时库中被识别；
    \item \textbf{片上网络流量模式}：通过硬件性能计数器（如 NoC 端口利用率、VC 拥塞统计）可间接推断是否存在基于门控结果的预测路由与预取策略。
    \end{enumerate}

\end{document}
